\documentclass[11pt]{article}
\usepackage[T1]{fontenc}
\usepackage[utf8]{inputenc}
\usepackage[margin=1in]{geometry}
\usepackage{amsmath,amssymb,amsthm}
\usepackage[hidelinks]{hyperref}

\newtheorem{theorem}{Theorem}
\newtheorem{lemma}{Lemma}
\newtheorem{corollary}{Corollary}
\newtheorem{definition}{Definition}
\newtheorem{remark}{Remark}
\newtheorem*{problem875}{Erd\H{o}s Problem \#875, in the asymptotic absolute-gap form used here}

\newcommand{\N}{\mathbb{N}}
\newcommand{\Z}{\mathbb{Z}}
\newcommand{\eps}{\varepsilon}
\newcommand{\D}{\Delta}
\newcommand{\Sig}{\Sigma}

\title{On an asymptotic absolute-gap form of Erd\H{o}s Problem \#875:\\
An infinite admissible set with gaps $o(n^{3+2\sqrt2})$}
\author{Draft note}
\date{May 2026}

\begin{document}
\maketitle

\begin{abstract}
Erd\H{o}s Problem \#875 asks how slowly an infinite admissible set can grow.  One associated absolute-gap question asks whether, after ignoring finitely many initial terms, the gaps can be bounded by a power of the index.  We give an explicit construction for this eventual absolute-gap question.  Namely, we construct an infinite admissible set $A=\{a_1<a_2<\cdots\}\subset\N$ such that
\[
        a_{n+1}-a_n=o\bigl(n^{3+2\sqrt2}\bigr).
\]
Consequently the eventual inequality $a_{n+1}-a_n\le n^{3+2\sqrt2}$ holds for all sufficiently large $n$.  This is a partial positive answer to the absolute polynomial-gap formulation; the exponent is not claimed to be optimal, and no claim is made about the relative-gap condition $a_{n+1}/a_n\to1$.  The construction is by blocks: the inter-block gap equals the internal spacing of the preceding block, although each block itself has large internal multiplicative spread.  The proof uses a carry invariant for cardinality-graded subset-sum differences.
\end{abstract}

\section{Introduction}

Throughout, $\N=\{1,2,3,\ldots\}$.  The symbols $\ll$, $\gg$, $\asymp$, $O(\cdot)$, $o(\cdot)$, and $\omega(\cdot)$ have their usual asymptotic meanings; all implied constants are absolute unless a dependence is explicitly indicated.

For a set $A\subset\N$ and an integer $r\ge1$, write
\[
        S_r(A)=\left\{\sum_{a\in U}a: U\subseteq A,\ |U|=r\right\}.
\]

\begin{problem875}
Let $A=\{a_1<a_2<\cdots\}\subset\N$ be an infinite set such that the sets $S_r(A)$ are disjoint for distinct $r\ge1$.  How fast can such a sequence grow?  How small can the gaps $a_{n+1}-a_n$ be?  In the asymptotic absolute-gap form considered here, for which real numbers $c$ is it possible, after ignoring finitely many initial terms, that
\[
        a_{n+1}-a_n\le n^c
\]
for all remaining $n$?
\end{problem875}

Throughout this note, every asserted polynomial gap bound is meant in this eventual sense.  We do not try to enforce a prescribed inequality for the finitely many initial values of $n$; doing so would require a separate finite adjustment.

The Erd\H{o}s Problems site records this as Erd\H{o}s Problem \#875, a problem of Deshouillers and Erd\H{o}s and an infinite version of Erd\H{o}s Problem \#874, the corresponding finite extremal problem \cite{Erdos1998,Bloom874,Bloom875}.  This note gives an explicit construction for the eventual absolute polynomial-gap part of Problem \#875, improving the exponent obtained from the Erd\H{o}s--Nicolas--S\'ark\"ozy construction by the simple counting implication described below.  We make no optimality claim for Problem \#875 as a whole, do not classify all possible exponents, do not address the relative-gap question, and do not attempt a complete historical survey beyond the sources cited here.

We use the following terminology, following the standard terminology in this problem.  A set $A\subset\N$ is called \emph{admissible} if, whenever $U,V\subseteq A$ are finite and
\[
        \sum_{u\in U}u=\sum_{v\in V}v,
\]
then $|U|=|V|$.  Equivalently, the sets $S_r(A)$ are pairwise disjoint.  Including the empty subset in the definition makes no difference for subsets of $\N$, since no nonempty subset has sum $0$.

The main result is the following partial positive answer to the eventual absolute-gap question in Erd\H{o}s Problem \#875.

\begin{theorem}[An admissible set with gaps $o(n^{3+2\sqrt2})$]\label{thm:main}
There is an infinite admissible set $A=\{a_1<a_2<\cdots\}\subset\N$ such that
\[
        a_{n+1}-a_n=o\bigl(n^{3+2\sqrt2}\bigr).
\]
Consequently,
\[
        a_{n+1}-a_n\le n^{3+2\sqrt2}
\]
for all sufficiently large $n$.
\end{theorem}

The theorem is intentionally not a literal small-$n$ statement.  In the construction below the first block is $B_1=\{1,4\}$, so the first gap is $3>1^{3+2\sqrt2}$.

Theorem~\ref{thm:main} concerns absolute gaps, not relative gaps.  With the notation of the construction below, the first two elements of the $j$th block are $M_j$ and $M_j(k_j+2)$, so their ratio is $k_j+2$.  Since $N_j\to\infty$ and $k_j=\lceil 2(1+N_{j-1})^{1+\sqrt2}\rceil$, these relative jumps tend to infinity.  On the other hand, the inter-block ratio is small:
\[
        \frac{\min B_{j+1}}{\max B_j}
        =\frac{M_{j+1}}{M_j k_j^2}
        =1+\frac1{k_j}+\frac1{k_j^2}
        \longrightarrow 1.
\]
Thus the construction separates the eventual absolute-gap issue from any claim about consecutive ratios tending to $1$.

For comparison with the finite theory, let $K(N)$ be the maximum possible cardinality of an admissible subset of $\{1,\ldots,N\}$.  This finite problem goes back to Erd\H{o}s and is recorded as Erd\H{o}s Problem \#874, which asks whether $K(N)\sim 2N^{1/2}$ \cite{Erdos1962,Bloom874}.  Straus introduced the term admissible and gave lower-bound constructions showing that the constant $2$ cannot be improved \cite{Straus1966}.  Deshouillers and Freiman proved the asymptotic upper bound
\[
        K(N)\le (2+o(1))\sqrt N
\]
\cite{DF1}, and later work of Deshouillers and Freiman settled the finite problem for all large $N$ \cite{DF2,Bloom874}.  Some related literature studies stronger stratified variants of finite subset-sum separation; here we use only the admissibility condition defined above \cite{Nicolas1999}.

The finite upper bound gives a simple necessary growth condition for any infinite admissible set.  If $A=\{a_1<a_2<\cdots\}$ is admissible, then $A\cap[1,a_n]$ is an admissible subset of $\{1,\ldots,a_n\}$ with at least $n$ elements.  Thus, for every fixed $\eps>0$ and all sufficiently large $n$,
\[
        n\le (2+\eps)\sqrt{a_n},
        \qquad\text{so}\qquad
        a_n\ge \frac{n^2}{(2+\eps)^2}.
\]
Equivalently,
\[
        a_n\ge (1/4-o(1))n^2.
\]
In particular, a gap bound $a_{n+1}-a_n\le n^c$ for all sufficiently large $n$ is impossible when $c<1$: if $c<0$, then $n^c<1$ eventually although the gaps are positive integers, while if $0\le c<1$, such a bound would imply $a_n=O(n^{c+1})=o(n^2)$.  This is only a necessary condition; it does not suggest that the borderline exponent $c=1$ is attainable.

Erd\H{o}s, Nicolas, and S\'ark\"ozy constructed an infinite admissible set with
\[
        |A\cap[1,x]|\gg x^{5-2\sqrt6}
\]
\cite{ENS1991}.  For that construction,
\[
        a_n\ll n^{1/(5-2\sqrt6)}=n^{5+2\sqrt6}.
\]
Consequently $a_{n+1}-a_n\le a_{n+1}\ll n^{5+2\sqrt6}$, and the constant implicit in $\ll$ can be absorbed into $n^{c-(5+2\sqrt6)}$ for every $c>5+2\sqrt6$.  Hence that construction already gives $a_{n+1}-a_n\le n^c$ eventually for every exponent strictly larger than $5+2\sqrt6$.  The construction below improves this ENS-derived absolute polynomial-gap comparison to the explicit exponent $3+2\sqrt2$ in the little-$o$ sense.  As recorded in Corollary~\ref{cor:growth}, it also gives
\[
        |A\cap[1,x]|=\omega\bigl(x^{1/(4+2\sqrt2)}\bigr),
\]
and hence, in particular, $|A\cap[1,x]|\gg x^\beta$ for every fixed $\beta<1/(4+2\sqrt2)$.  Here $1/(4+2\sqrt2)\approx0.1464$, compared with $5-2\sqrt6\approx0.1010$ in the cited Erd\H{o}s--Nicolas--S\'ark\"ozy construction.

\section{Graded differences}

For a finite set $B\subset\Z$, write
\[
        \Sig(B)=\sum_{b\in B}b.
\]
When $B\subset\N$, every subset sum from $B$ lies in $[0,\Sig(B)]$; this positivity is the only context in which we use $\Sig(B)$ as an upper bound for subset sums.

For a finite set $B\subset\Z$ and an integer $t$, define
\[
\D_t(B)=
\left\{
\sum_{u\in U}u-\sum_{v\in V}v:
U,V\subseteq B,\ |U|-|V|=t
\right\}.
\]
Empty subsets are allowed in this definition.  This convention is harmless for positive sets and is part of the definition for general finite subsets of $\Z$.  Thus $0\in \D_t(B)$ means that two subsets of $B$ of cardinality difference $t$ have the same sum.  For finite sets of integers, possibly including negative integers, we use the same word admissible for this cardinality-graded subset-sum condition: $B$ is admissible if no equality of two subset sums occurs with different subset cardinalities.

\begin{lemma}\label{lem:admissible-delta}
A finite set $B\subset\Z$ is admissible if and only if
\[
        0\notin \D_t(B)\qquad\text{for every integer }t\ne 0.
\]
\end{lemma}

\begin{proof}
This is just the definition.  An equality of two subset sums with cardinality difference $t$ is exactly an element $0\in \D_t(B)$.
\end{proof}

For an integer $k\ge1$, put
\[
        C_k=\{1+i(k+1):0\le i<k\}.
\]
Every element of $C_k$ is congruent to $1$ modulo $k+1$.  Therefore any sum of $r$ distinct elements of $C_k$ is congruent to $r$ modulo $k+1$.  Since $0\le r\le k$, this residue determines $r$.  Thus, if two subset sums from $C_k$ are equal, the two subset cardinalities are equal, and hence $C_k$ is admissible.  We shall use the elementary facts
\[
        |C_k|=k,
        \qquad
        \max C_k=k^2,
        \qquad
        \Sig(C_k)=\frac{k^3+k}{2}.
\]
Also define
\[
        Q_k=k^2+k+1.
\]
For an integer $M$ and a finite set $B\subset\Z$, write
\[
        MB=\{Mb:b\in B\}
\]
for the dilation of $B$ by $M$.

\section{The local carry lemma}

The key point is that a small error near a multiple of $Q_k$ determines the cardinality defect.

\begin{lemma}[local carry]\label{lem:local-carry}
Let $m\ge1$ be an integer, and put
\[
        C_m=\{1+i(m+1):0\le i<m\},
        \qquad
        Q=m^2+m+1.
\]
Suppose $Y\in \D_t(C_m)$ and
\[
        Y=Qw-u
\]
with $u,w\in\Z$ satisfying
\[
        |u|<\frac m4,
        \qquad
        |w|<\frac m2.
\]
Then
\[
        t=w-u.
\]
\end{lemma}

\begin{proof}
Choose subsets realizing the given membership $Y\in\D_t(C_m)$.  After cancelling their common part, write
\[
        Y=\sum_{i=0}^{m-1}\eps_i\bigl(1+i(m+1)\bigr),
        \qquad
        \eps_i\in\{-1,0,1\},
\]
where $\eps_i=1$ if the $i$th element of $C_m$ lies only in the first subset, $\eps_i=-1$ if it lies only in the second subset, and $\eps_i=0$ otherwise.  Cancelling common elements changes neither the difference $Y$ nor the cardinality defect $t$.  Then
\[
        t=\sum_i\eps_i,
        \qquad
        Y=t+(m+1)I,
        \qquad
        I=\sum_i i\eps_i.
\]
Since $Q=m^2+m+1\equiv1\pmod {m+1}$, the equation $Y=Qw-u$ gives
\[
        t\equiv w-u\pmod {m+1}.
\]
Put $t_0=w-u$.  From $|w|<m/2$ and $|u|<m/4$ we have $|t_0|<3m/4$.  Also $|t|\le m$ from the representation $t=\sum_i\eps_i$ with $m$ summands in $\{-1,0,1\}$.  If $t=t_0+\ell(m+1)$ with $|\ell|\ge2$, then
\[
        |t|\ge 2(m+1)-|t_0|>m,
\]
which is impossible.  Hence the only possible values of $t$ are
\[
        t=t_0,
        \qquad
        t=t_0+(m+1),
        \qquad
        t=t_0-(m+1).
\]
The two shifted cases are the only possible wraparound obstructions modulo $m+1$.  We rule them out by extremal bounds on $I$.

Assume first that
\[
        t=t_0+(m+1).
\]
Let $a=u-w=-t_0$.  Then
\[
        t=m+1-a.
\]
Since $t=\sum_i\eps_i\le m$, we have $a\ge1$.  The identity $t+(m+1)I=Qw-u$, together with $u=w+a$, gives
\[
        m+1-a+(m+1)I=(m^2+m+1)w-(w+a)=(m+1)mw-a,
\]
so
\[
        I=mw-1.
\]
On the other hand, define $d_i=1-\eps_i\in\{0,1,2\}$.  Since
\[
        \sum_i d_i=m-t=a-1,
\]
the total $a-1$ is nonnegative, and since $0\le i\le m-1$ and $d_i\ge0$, we have
\[
        \sum_i i d_i\le (m-1)\sum_i d_i=(a-1)(m-1).
\]
Therefore
\[
        I
        =\frac{m(m-1)}2-\sum_i i d_i
        \ge
        L:=\frac{m(m-1)}2-(a-1)(m-1).
\]
But $u=w+a$ and $|u|<m/4$, so $w<m/4-a$.  Hence
\[
        I=mw-1<R:=\frac{m^2}{4}-ma-1.
\]
These two inequalities are incompatible, since
\[
        L-R
        =
        \left(
        \frac{m(m-1)}2-(a-1)(m-1)
        \right)
        -
        \left(
        \frac{m^2}{4}-ma-1
        \right)
        =
        \frac{m^2}{4}+\frac m2+a>0.
\]
Thus we would have $I\ge L>R>I$, a contradiction.  Hence $t=t_0+(m+1)$ is impossible.

If $t=t_0-(m+1)$, then $-Y\in\D_{-t}(C_m)$ and
\[
        -Y=Q(-w)-(-u).
\]
The hypotheses $|u|<m/4$ and $|w|<m/2$ are preserved under this sign change.  For the negated difference the corresponding residue representative is
\[
        (-w)-(-u)=-t_0,
\]
and the assumed relation gives
\[
        -t=-t_0+(m+1).
\]
This is exactly the already excluded positive shifted case with $m$ unchanged and with $(Y,t,u,w)$ replaced by $(-Y,-t,-u,-w)$.  Therefore the two wraparound alternatives are impossible, leaving the unique residue-compatible possibility
\[
        t=t_0=w-u.
\]
\end{proof}

\section{Carry-good pairs and closure}

The next invariant is the mechanism that lets the construction place a new block close to the old scale.  Informally, if a cardinality-graded difference from the old set lands exactly on a multiple of the current modulus $M$, then the multiple records the cardinality defect.  This is the ``carry'' that will be passed to the local block.

\begin{definition}
Let $A\subset\N$ be finite and let $M$ be a positive integer.  We say that $(A,M)$ is \emph{carry-good} if:
\begin{enumerate}
    \item $A\subset [1,M)\cap\N$;
    \item whenever $D\in \D_T(A)$ and $D=Mq$ for some $q\in\Z$, one has $T=q$.
\end{enumerate}
\end{definition}

By Lemma~\ref{lem:admissible-delta}, a carry-good pair automatically gives an admissible finite positive set: if $0\in \D_T(A)$, then $0=M\cdot0$, so the carry-good condition gives $T=0$.

\begin{remark}[Starting after a finite prefix]\label{rem:prefix}
Finite initial gaps are irrelevant for the eventual form of Problem~\#875.  If $A_0\subset\N$ is any fixed finite admissible set, choose an initial modulus $M>\Sig(A_0)$.  Then $A_0\subset[1,M)\cap\N$.  If $D\in\D_T(A_0)$ and $D=Mq$, then the two subset sums defining $D$ lie in $[0,\Sig(A_0)]$, so $|D|<M$ and hence $q=0$.  Since $A_0$ is admissible, $D=0$ forces $T=0$.  Thus $(A_0,M)$ is carry-good.

The first application of the closure lemma is also legitimate.  If $A_0=\varnothing$, then $\Sig(A_0)/M=0$.  If $A_0\ne\varnothing$ and $N_0=|A_0|$, then $\Sig(A_0)/M<1$, while the first new block length in the construction would be
\[
        k_1=\left\lceil 2(1+N_0)^{1+\sqrt2}\right\rceil\ge 11,
\]
so $\Sig(A_0)/M<1<k_1/4$.  Hence one may run the same construction with this pair as the initial pair, taking the initial cardinality to be $N_0=|A_0|$ and the initial modulus to be $M$.  This remark preserves eventual gap estimates only; it does not by itself impose any desired inequalities on the finitely many initial gaps.
\end{remark}

\begin{lemma}[carry closure]\label{lem:closure}
Suppose $(A,M)$ is carry-good, and let $k\ge1$ be an integer such that
\[
        \frac{\Sig(A)}M<\frac k4.
\]
Define
\[
        A'=A\cup M C_k,
        \qquad
        M'=M Q_k=M(k^2+k+1).
\]
Then $(A',M')$ is carry-good.  Moreover,
\[
        \frac{\Sig(A')}{M'}<\frac k2.
\]
\end{lemma}

\begin{proof}
The strict inequality $\Sig(A)/M<k/4$ is used to satisfy the strict small-carry hypothesis in Lemma~\ref{lem:local-carry}; the strictness in the conclusion will similarly be used at the next stage.  Because $A\subset[1,M)\cap\N$ and every element of $MC_k$ is at least $M$, the union defining $A'$ is disjoint.  Also every element of $A'$ is smaller than $M'$.  Indeed, every element of $A$ is smaller than $M<M'$, and
\[
        \max(MC_k)=Mk^2<M(k^2+k+1)=M'.
\]
Thus $A'\subset[1,M')\cap\N$.

Let
\[
        D\in\D_T(A'),
        \qquad
        D=M'q.
\]
Choose subsets $U,V\subseteq A'$ that realize $D$, and split each of them uniquely into its old part and its new part:
\[
        U=U_A\sqcup MU_C,
        \qquad
        V=V_A\sqcup MV_C,
\]
where $U_A,V_A\subseteq A$ and $U_C,V_C\subseteq C_k$.  Define the old and new cardinality defects by
\[
        T_A=|U_A|-|V_A|,
        \qquad
        t=|U_C|-|V_C|.
\]
Then $T=T_A+t$.  The new contribution is $MY$, where
\[
        Y=\sum_{u\in U_C}u-\sum_{v\in V_C}v\in\D_t(C_k).
\]
Thus
\[
        D=D_A+MY,
\]
where
\[
        D_A\in \D_{T_A}(A),
        \qquad
        Y\in \D_t(C_k),
        \qquad
        T=T_A+t.
\]
Since $M'=M Q_k$, the equation $D=M'q$ gives
\[
        D_A+MY=M Q_k q.
\]
Hence
\[
        D_A=M(Q_k q-Y)=Mu
\]
for the integer $u=Q_k q-Y$.  Both subset sums defining $D_A$ lie in $[0,\Sig(A)]$, so $|D_A|\le \Sig(A)$; consequently
\[
        |u|=\frac{|D_A|}{M}\le \frac{\Sig(A)}M<\frac k4.
\]
Because $(A,M)$ is carry-good, we get
\[
        T_A=u.
\]
The equation then becomes
\[
        Y=Q_k q-u.
\]
It remains to verify the size condition on $q$ needed for Lemma~\ref{lem:local-carry}.  Since $D=M'q$ is a difference of two subset sums from the positive set $A'$, both subset sums lie in $[0,\Sig(A')]$, and so
\[
        |q|\le \frac{\Sig(A')}{M'}.
\]
Using disjointness of the union defining $A'$, we compute
\[
\frac{\Sig(A')}{M'}
=
\frac{\Sig(A)}{M Q_k}+\frac{\Sig(C_k)}{Q_k}
=
\frac{\Sig(A)}{M Q_k}+\frac{k^3+k}{2(k^2+k+1)}.
\]
The first term is $<k/(4Q_k)$.  The second term is below $k/2$ by
\[
        \frac{k}{2}-\frac{k^3+k}{2Q_k}=\frac{k^2}{2Q_k}.
\]
Since $k/(4Q_k)<k^2/(2Q_k)$ for $k\ge1$, it follows that
\[
        \frac{\Sig(A')}{M'}<\frac k2.
\]
This estimate is also the quantity that will be compared with the next block length in the induction.  In the present closure step, it gives $|q|<k/2$.  Applying Lemma~\ref{lem:local-carry} with local carry $q$ gives
\[
        t=q-u.
\]
Consequently
\[
        T=T_A+t=u+(q-u)=q.
\]
Thus $(A',M')$ is carry-good, and the displayed estimate for $\Sig(A')/M'$ has also been proved.
\end{proof}

\section{The infinite construction and the gap estimate}

\begin{proof}[Proof of Theorem~\ref{thm:main}]
Start with
\[
        A_0=\varnothing,
        \qquad
        M_1=1,
        \qquad
        N_0=0.
\]
The pair $(A_0,M_1)$ is carry-good: $A_0\subset[1,M_1)\cap\N$, and the only graded difference from $A_0$ is $D=0$ with cardinality defect $T=0$.

Put
\[
        \theta=1+\sqrt2,
        \qquad
        \tau=2+\sqrt2,
        \qquad
        f=3+2\sqrt2.
\]
These parameters satisfy
\[
        f=\theta+\tau,
        \qquad
        \tau=\theta\sqrt2,
        \qquad
        \tau+\theta(2-\tau)=0.
\]

For $j\ge1$, define recursively
\[
        k_j=\left\lceil 2(1+N_{j-1})^{\theta}\right\rceil,
        \qquad
        B_j=M_j C_{k_j},
        \qquad
        A_j=A_{j-1}\cup B_j,
\]
\[
        M_{j+1}=M_j Q_{k_j},
        \qquad
        N_j=|A_j|.
\]
For instance, $k_1=2$ and $B_1=C_2=\{1,4\}$.  The next values are $M_2=7$, $N_1=2$, and $k_2=29$, so every gap whose left endpoint lies in $B_2$ is $7\cdot 30=210$.

We prove the following induction claim.  At stage $j$, the old set $A_{j-1}$ is controlled modulo $M_j$, and after adjoining $B_j$ the new set is controlled modulo $M_{j+1}$.  More precisely, for every $j\ge1$, the block $B_j$ is disjoint from $A_{j-1}$,
\[
        M_{j+1}\in\N,
        \qquad
        A_j\subset[1,M_{j+1})\cap\N,
        \qquad
        N_j=N_{j-1}+k_j,
\]
$(A_j,M_{j+1})$ is carry-good, and
\[
        \frac{\Sig(A_j)}{M_{j+1}}<\frac{k_j}{2}.
\]
For $j=1$, $M_2=M_1Q_{k_1}$ is a positive integer, the block is disjoint from $A_0$, and $N_1=N_0+k_1$.  The closure lemma applies to $(A_0,M_1)$ because $\Sig(A_0)/M_1=0<k_1/4$, and it gives the remaining base-case claims.

Now suppose the claims have been proved through stage $j-1$, with $j\ge2$.  Then $M_{j+1}=M_j Q_{k_j}$ is a positive integer.  Also $k_{j-1}\le N_{j-1}$ because $N_{j-1}=k_1+\cdots+k_{j-1}$, and
\[
        k_j\ge 2(1+N_{j-1})^\theta>2N_{j-1},
\]
since $\theta>1$.  Hence the previous stage gives
\[
        \frac{\Sig(A_{j-1})}{M_j}<\frac{k_{j-1}}2\le \frac{N_{j-1}}2
        <\frac{k_j}{4}.
\]
This strict final inequality is one role of the factor $2$ in the definition of $k_j$; the same normalization also leaves the geometric contraction constant below strictly smaller than $1$.  Lemma~\ref{lem:closure} applies to $(A_{j-1},M_j)$ with parameter $k_j$.  It gives that $(A_j,M_{j+1})$ is carry-good and that
\[
        \frac{\Sig(A_j)}{M_{j+1}}<\frac{k_j}{2}.
\]
The same lemma also gives $A_j\subset[1,M_{j+1})\cap\N$.  Since $A_{j-1}\subset[1,M_j)\cap\N$ and the first element of $B_j=M_jC_{k_j}$ is $M_j$, the new block is disjoint from all earlier elements, so $N_j=N_{j-1}+k_j$.  This completes the induction.

Set
\[
        A=\bigcup_{j\ge1} A_j=\bigsqcup_{j\ge1}B_j.
\]
Each $k_j\ge2$, and each stage adds a new disjoint block, so $A$ is infinite.  The set $A$ is admissible: if finite subsets $U,V\subseteq A$ have the same sum, then $U\cup V\subseteq A_J$ for some $J$, and the finite set $A_J$ is admissible because $(A_J,M_{J+1})$ is carry-good.  Therefore $|U|=|V|$.

It remains to estimate the gaps in the increasing enumeration of $A$.  The induction above shows that adjacent blocks occur in increasing order.  Indeed,
\[
        \max B_j=M_j k_j^2<M_j(k_j^2+k_j+1)=M_{j+1}=\min B_{j+1}.
\]
By induction over adjacent blocks, every element of $B_i$ is smaller than every element of $B_j$ whenever $i<j$.  Thus the increasing enumeration of $A$ is obtained by listing the blocks $B_j$ in order.  Also $N_j\to\infty$, because $N_j=N_{j-1}+k_j$ and $k_j\ge2$.

For $j\ge2$, define
\[
        P_j=\frac{M_j}{N_{j-1}^{\tau}}.
\]
Since $k_j\ge2$, we have
\[
        Q_{k_j}=k_j^2+k_j+1\le2k_j^2.
\]
Since $N_j\ge k_j$ and $\tau>0$, it follows that
\[
        P_{j+1}
        =\frac{M_{j+1}}{N_j^\tau}
        \le
        \frac{2M_j k_j^2}{k_j^\tau}
        =2M_j k_j^{2-\tau}.
\]
Here $2-\tau=-\sqrt2<0$.  Since the exponent is negative, the lower bound $k_j\ge2N_{j-1}^{\theta}$ gives the upper bound
\[
        k_j^{2-\tau}\le \bigl(2N_{j-1}^{\theta}\bigr)^{2-\tau}.
\]
Therefore, for $j\ge2$,
\[
        P_{j+1}
        \le
        2M_j\bigl(2N_{j-1}^{\theta}\bigr)^{-\sqrt2}
        =2^{1-\sqrt2}P_jN_{j-1}^{\tau-\theta\sqrt2}
        =2^{1-\sqrt2}P_j.
\]
Since $P_2$ is finite and $2^{1-\sqrt2}<1$, it follows that $P_j\to0$ geometrically.

We now estimate all gaps whose left endpoint lies in a block $j\ge2$; the first block, including the gap from $B_1$ to $B_2$, contributes only finitely many gaps.  Inside the block $B_j$, the elements are
\[
        M_j\bigl(1+i(k_j+1)\bigr),
        \qquad 0\le i<k_j,
\]
so each consecutive internal gap is
\[
        M_j(k_j+1).
\]
The gap between the last element of $B_j$ and the first element of $B_{j+1}$ is
\[
        M_{j+1}-M_j k_j^2=M_j(k_j+1).
\]
Hence every gap whose left endpoint lies in the $j$th block is exactly $M_j(k_j+1)$.

For $j\ge2$ we have $N_{j-1}\ge2$, and hence there is an absolute constant $C_0$ such that
\[
        k_j+1
        \le C_0N_{j-1}^{\theta}.
\]
Indeed, this follows immediately from
\[
        k_j+1
        \le 2(1+N_{j-1})^\theta+2
        \le \bigl(2(3/2)^\theta+1\bigr)N_{j-1}^\theta.
\]

For an index $n$ whose gap is not among the finitely many gaps with left endpoint in $B_1$, let $j(n)$ be the unique block index such that $a_n\in B_{j(n)}$.  Then $j(n)\to\infty$ as $n\to\infty$, because any finite union of blocks contains only finitely many elements.  If $j=j(n)\ge2$, then $N_{j-1}<n\le N_j$ and therefore
\[
        \frac{a_{n+1}-a_n}{n^f}
        =\frac{M_j(k_j+1)}{n^f}
        \le
        \frac{M_j(k_j+1)}{N_{j-1}^{\tau+\theta}}
        \le C_0\frac{M_j}{N_{j-1}^{\tau}}
        =C_0P_j,
\]
since $f=\tau+\theta$.  Equivalently,
\[
        \sup_{N_{j-1}<n\le N_j}
        \frac{a_{n+1}-a_n}{n^f}
        \le C_0P_j
        \qquad (j\ge2).
\]
Because $P_j\to0$ and $j(n)\to\infty$, we obtain
\[
        a_{n+1}-a_n=o(n^f),
        \qquad
        f=3+2\sqrt2.
\]
The eventual inequality $a_{n+1}-a_n\le n^{3+2\sqrt2}$ follows because the ratio tends to $0$, and is therefore eventually at most $1$.
\end{proof}

\begin{remark}[Choice of parameters]\label{rem:parameters}
The following heuristic explains the parameter choice in the present proof; it is not an optimality claim for Problem~\#875.  If one chooses block lengths of size $k_j\asymp N_{j-1}^{\theta}$ and tries to prove $M_j=o(N_{j-1}^{\tau})$ using the recurrence above, then the argument requires
\[
        \tau+\theta(2-\tau)\le0,
        \qquad \tau>2.
\]
Thus $\theta\ge \tau/(\tau-2)$, and minimizing $\theta+\tau$ gives $\tau=2+\sqrt2$ and $\theta=1+\sqrt2$.  For the concrete normalization used above, including the coefficient $2$ in the definition of $k_j$, the remaining constant is $2^{1-\sqrt2}<1$, so equality in the exponent condition still gives geometric decay.
\end{remark}

\begin{corollary}\label{cor:growth}
The set constructed in Theorem~\ref{thm:main} satisfies the following simple consequence of the gap estimate:
\[
        a_n=o\bigl(n^{4+2\sqrt2}\bigr),
        \qquad
        |A\cap[1,x]|=\omega\bigl(x^{1/(4+2\sqrt2)}\bigr).
\]
\end{corollary}

\begin{proof}
Let $f=3+2\sqrt2$.  From the theorem,
\[
        a_{r+1}-a_r=o(r^f).
\]
We show that summing preserves the little-$o$ estimate.  Given $\varepsilon>0$, choose an integer $R\ge1$ such that
\[
        a_{r+1}-a_r\le \varepsilon r^f
\]
for all $r\ge R$.  Then, for $n>R$,
\[
        a_n
        =a_R+\sum_{r=R}^{n-1}(a_{r+1}-a_r)
        \le
        a_R+\varepsilon\sum_{r=R}^{n-1}r^f
        \le
        a_R+C_f\varepsilon n^{f+1},
\]
where $C_f$ depends only on $f$.  Dividing by $n^{f+1}$ and taking the limsup gives
\[
        \limsup_{n\to\infty}\frac{a_n}{n^{f+1}}\le C_f\varepsilon.
\]
Since $\varepsilon$ is arbitrary,
\[
        a_n=o(n^{f+1})=o\bigl(n^{4+2\sqrt2}\bigr).
\]

Put $\alpha=1/(f+1)=1/(4+2\sqrt2)$.  To prove the counting estimate, fix any constant $L>0$ and set $n=\lfloor Lx^\alpha\rfloor$.  For all sufficiently large $x$, this integer satisfies $n\ge1$.  As $x\to\infty$, we have $n\to\infty$ and
\[
        \frac{a_n}{x}
        =
        \frac{a_n}{n^{f+1}}\cdot\frac{n^{f+1}}{x}
        \longrightarrow 0,
\]
because $n^{f+1}/x$ remains bounded for fixed $L$.  Hence $a_n\le x$ for all sufficiently large $x$, and therefore
\[
        |A\cap[1,x]|\ge n\ge \frac L2 x^\alpha
\]
for all sufficiently large $x$.  Since $L$ was arbitrary, $|A\cap[1,x]|=\omega(x^\alpha)$.
\end{proof}

\begin{remark}[Possible improvements]\label{rem:improvements}
The local block $C_k$ has size $k$ and maximum $k^2$.  It is therefore not finite-cardinality optimal in the constant-factor sense, since finite admissible subsets of $[1,N]$ can have about $2\sqrt N$ elements.  Within the present block-carry scheme, finite-cardinality optimality alone would not be enough to improve the exponent.  A replacement block would also need a compatible carry lemma, a controlled first element, a controlled maximum, small internal consecutive gaps, and a small edge gap between the maximum of the block and the next modulus.  In the present block $C_k$, this edge gap is
\[
        Q_k-\max C_k=k+1,
\]
which exactly matches the internal spacing.  Thus improving the exponent by this route would require more than substituting an arbitrary finite extremal admissible set; it appears to require a more radical carry-compatible local structure.
\end{remark}

\begin{thebibliography}{10}

\bibitem{Erdos1962}
P. Erd\H{o}s,
\emph{Sz\'amelm\'eleti megjegyz\'esek, III.\ N\'eh\'any addit\'iv sz\'amelm\'eleti probl\'em\'ar\'ol},
Matematikai Lapok \textbf{13} (1962), 28--38.

\bibitem{Erdos1998}
P. Erd\H{o}s,
\emph{Some of my new and almost new problems and results in combinatorial number theory},
in \emph{Number Theory: Diophantine, Computational and Algebraic Aspects},
Proceedings of the International Conference held in Eger, Hungary, July 29--August 2, 1996,
edited by K. Gy\H{o}ry, A. Peth\H{o}, and V. T. S\'os,
De Gruyter, Berlin/New York, 1998, pp. 169--180.
DOI: \href{https://doi.org/10.1515/9783110809794.169}{10.1515/9783110809794.169}.

\bibitem{Straus1966}
E. G. Straus,
\emph{On a problem in combinatorial number theory},
Journal of Mathematical Sciences \textbf{1} (1966), 77--80.

\bibitem{ENS1991}
P. Erd\H{o}s, J.-L. Nicolas, and A. S\'ark\"ozy,
\emph{Sommes de sous-ensembles},
S\'eminaire de th\'eorie des nombres de Bordeaux, S\'erie 2, \textbf{3} (1991), no. 1, 55--72.
DOI: \href{https://doi.org/10.5802/jtnb.42}{10.5802/jtnb.42}.

\bibitem{DF1}
J.-M. Deshouillers and G. A. Freiman,
\emph{On an additive problem of Erd\H{o}s and Straus, 1},
Israel Journal of Mathematics \textbf{92} (1995), no. 1--3, 33--43.
DOI: \href{https://doi.org/10.1007/BF02762069}{10.1007/BF02762069}.

\bibitem{DF2}
J.-M. Deshouillers and G. A. Freiman,
\emph{On an additive problem of Erd\H{o}s and Straus, 2},
in \emph{Structure theory of set addition}, Ast\'erisque no. \textbf{258} (1999), 141--148.
Available at \url{https://www.numdam.org/item/AST_1999__258__141_0/}.

\bibitem{Nicolas1999}
J.-L. Nicolas,
\emph{Stratified sets},
in \emph{Structure theory of set addition}, Ast\'erisque no. \textbf{258} (1999), 205--215.
Available at \url{https://www.numdam.org/item/AST_1999__258__205_0/}.

\bibitem{Bloom874}
T. F. Bloom,
\emph{Erd\H{o}s Problem \#874},
\url{https://www.erdosproblems.com/874}, accessed May 7, 2026.

\bibitem{Bloom875}
T. F. Bloom,
\emph{Erd\H{o}s Problem \#875},
\url{https://www.erdosproblems.com/875}, accessed May 7, 2026.

\end{thebibliography}

\end{document}
